\documentclass[11pt]{amsart}
\usepackage{amsmath,amssymb,amsthm}
\usepackage{booktabs}
\usepackage[hidelinks]{hyperref}
\hypersetup{pdftitle={Covering codes checked by the Lean kernel: K7(9,4) <= 1351 and K2(6,1) = 12},pdfauthor={Thiago Patzdorf}}
\pdfstringdefDisableCommands{\def\\{ }\def\le{<=}}
\usepackage[margin=1in]{geometry}

\newtheorem{theorem}{Theorem}
\newtheorem{lemma}[theorem]{Lemma}
\theoremstyle{definition}
\newtheorem{definition}[theorem]{Definition}
\theoremstyle{remark}
\newtheorem*{remark}{Remark}

\title[Covering codes checked by the Lean kernel]{Covering codes checked by the Lean kernel:\\ $K_7(9,4)\le 1351$ and $K_2(6,1)=12$}
\author{Thiago Patzdorf}
\address{Universidade de Caxias do Sul, Brazil}
\email{thiagosleman@gmail.com}
\date{2 October 2026}

\begin{document}
\begin{abstract}
We report three results on the covering number $K_q(n,R)$, each proved in Lean~4 against Mathlib~v4.34.1 and checked by the kernel, with no \texttt{sorry} and no \texttt{native\_decide}.
First, an explicit code of $1351$ words covers $(\mathbb{Z}/7\mathbb{Z})^9$ at Hamming radius $4$, so $K_7(9,4)\le 1351$; the best upper bound we found in the literature is $1475$.
Second, $K_2(6,1)\ge 11$ follows from a short double-counting argument with no search.
Third, $K_2(6,1)=12$, a classical value, is proved by a pruned search whose soundness is itself a theorem, split into $38$ kernel decisions.
The same library contains a formal proof of the sphere-covering bound and eight numerical instances of it; those instances are below the best known lower bounds and are kept only as a test of the arithmetic.
Seven further codes improve the tables we consulted but are verified by computer only.
\end{abstract}
\maketitle

\section{Introduction}
The covering number $K_q(n,R)$ is the least cardinality of a code $C\subseteq A^n$, with $|A|=q$, such that every word of $A^n$ lies at Hamming distance at most $R$ from some word of $C$~\cite{cohen1997}.
Upper bounds come from constructions and lower bounds from counting; the tables of Kéri~\cite{keri} and the recent work of Gijswijt and Polak~\cite{gp2025} and of Marosi~\cite{marosi2026} record the current state.

Upper bounds are certificates: a list of words and the claim that it covers.
Checking the claim is a finite computation, but for $q^n$ in the tens of millions it is done by programs that are rarely audited.
Florath~\cite{florath2026} built a Lean library of covering-code bounds whose derivations are replayed by the kernel; it formalizes the sphere-covering bound and records many cells.
This note adds a certificate that the kernel checks word by word for a cell where the code is new, and a formal determination of a small classical value.

The proofs use Lean~4~\cite{lean4} and Mathlib~\cite{mathlib}. The source is \url{https://github.com/thiagopatzdorf/Matematica}, tag \texttt{v0.\allowbreak 3.\allowbreak 0}.
A statement below is a claim only when it is the name of a declaration in that tree.
The default target \texttt{CoveringLean} builds in seconds; the two search certificates form the target \texttt{CoveringHeavy}.

\section{Definitions}
In the development \texttt{CoveringA2} a word is a function $\mathtt{Fin}\,n\to\mathtt{ZMod}\,q$, distance is Mathlib's \texttt{hammingDist}, and
\[
  \texttt{Covers}\;R\;C \;:\iff\; \forall x,\ \exists c\in C,\ \mathrm{d}(x,c)\le R .
\]
In \texttt{CoveringA6} the alphabet is $\mathtt{Fin}\,q$, with the same distance and the same notion of covering, and
\[
  \texttt{IsK}\;q\;n\;R\;k \;:\iff\; \bigl(\exists C,\ |C|=k \wedge \texttt{Covers}\;R\;C\bigr) \wedge \forall C,\ |C|<k \to \neg\,\texttt{Covers}\;R\;C ,
\]
which is the equality $K_q(n,R)=k$.

\section{An upper bound: $K_7(9,4)\le 1351$}

\begin{theorem}\label{thm:k794}
There is $C\subseteq(\mathbb{Z}/7\mathbb{Z})^9$ with $|C|=1351$ that covers at radius $4$.
In Lean: \texttt{CoveringKernel.\allowbreak K7\_\allowbreak 9\_\allowbreak 4\_\allowbreak le\_\allowbreak 1351\_\allowbreak kernel}, stated as
\[
  \exists\, C : \mathtt{Finset}\,(\mathtt{Fin}\,9\to\mathtt{ZMod}\,7),\quad C.\mathtt{card}=1351 \;\wedge\; \texttt{Covers}\;4\;C .
\]
\end{theorem}

\subsection*{The code}
Words are encoded as natural numbers $\sum_i w_i 7^i$ and listed in strictly increasing order in \texttt{C1\_\allowbreak Data\_\allowbreak K7\_\allowbreak 9\_\allowbreak 4}; the same list is the text file \texttt{data/\allowbreak codes/\allowbreak q7\_\allowbreak n9\_\allowbreak R4\_\allowbreak M1351.\allowbreak txt}, whose canonical SHA-256 begins \texttt{54dbdade}.
The code is three cosets of a $[9,3]_7$ linear code ($1029$ words) together with $322$ words added greedily to cover what the cosets miss.
It was found with the generator \texttt{lincov} of Marosi's public repository \texttt{Mapika/coldcase}, which accompanies~\cite{marosi2026}; the program that added the last $322$ words was not recorded.
The theorem does not depend on how the code was found.

\subsection*{The proof}
A Boolean function \texttt{go} walks the coordinates one at a time, keeping for each prefix the codewords still within reach and their remaining budget of mismatches; \texttt{go\_\allowbreak sound} proves that \texttt{go} returning \texttt{true} implies that every word has a codeword at distance at most $R$.
\texttt{covers\_\allowbreak of\_\allowbreak go} turns that into \texttt{Covers}, using that Mathlib's distance agrees with the digit-wise count, and \texttt{cert\_\allowbreak of\_\allowbreak go} adds the cardinality, from a kernel check that every entry is below $7^9$ and the list is strictly increasing.
\texttt{go\_\allowbreak split} lets the first four coordinates be fixed in advance, so the computation becomes $7^4=2401$ leaf theorems, each closed by \texttt{decide +kernel}, grouped in $95$ files that the build compiles in parallel.
The same pipeline accepts the binary Hamming code of length $7$ and rejects a copy of it with one word replaced; both are tests in the library.

\begin{remark}
The best upper bound we found is $K_7(9,4)\le 1475$, by Marosi~\cite{marosi2026} (arXiv version 3, September 2026); Kéri's tables~\cite{keri} give $1843$.
The best lower bound we know is $264$.
That the bound is new is our reading of those sources on 1 October 2026, not something the kernel checks.
\end{remark}

\section{A lower bound without search: $K_2(6,1)\ge 11$}

\begin{lemma}\label{lem:even}
For distinct $a,c\in(\mathbb{F}_2)^6$, the balls of radius $1$ about $a$ and $c$ meet in $0$ or $2$ words.
\end{lemma}
\begin{proof}
At distance $1$ the intersection is $\{a,c\}$; at distance $2$ it is the two words between them; at distance at least $3$ it is empty.
In Lean the case $a=0$ is a kernel decision over $64$ words (\texttt{even\_\allowbreak inter\_\allowbreak zero}), and translation gives the general case (\texttt{inter\_\allowbreak card\_\allowbreak translate}).
\end{proof}

\begin{theorem}\label{thm:ge11}
Every code that covers $(\mathbb{F}_2)^6$ at radius $1$ has at least $11$ words (\texttt{CoveringA6.\allowbreak K\_\allowbreak 2\_\allowbreak 6\_\allowbreak 1\_\allowbreak ge\_\allowbreak 11}).
\end{theorem}
\begin{proof}
Let $K=|C|$ and let $m(y)$ be the number of codewords whose ball contains $y$, so $m(y)\ge 1$ and $\sum_y m(y)=7K$.
For each $x$ put $S(x)=\sum_{y\in B(x)} m(y)=\sum_{c\in C}|B(x)\cap B(c)|$.
Then $S(x)\ge 7$, and if $x\notin C$ every term is even by Lemma~\ref{lem:even}, so $S(x)\ge 8$.
Summing $S(x)+[x\in C]\ge 8$ over all $64$ words and using $\sum_x S(x)=7\sum_y m(y)=49K$ gives $50K\ge 512$, hence $K\ge 11$.
The counting is \texttt{excess\_\allowbreak bound}, stated for any family of symmetric balls of size $7$ with even pairwise intersections.
\end{proof}

This replaces a hypothesis that the previous version of this note left open: the statement \texttt{no\_\allowbreak cover\_\allowbreak 2\_\allowbreak 6\_\allowbreak 1\_\allowbreak le10\_\allowbreak uncond} is the earlier conditional one with the condition removed.

\section{An exact value: $K_2(6,1)=12$}

\begin{theorem}\label{thm:eq12}
\texttt{SC.\allowbreak K\_\allowbreak 2\_\allowbreak 6\_\allowbreak 1\_\allowbreak eq12} : \texttt{IsK}\;2\;6\;1\;12.
\end{theorem}

The value is classical~\cite{sk1968}.
The upper bound is an explicit $12$-word code (\texttt{code12}, encodings $0,1,2,15,28,23,57,58,59,39,44,52$).
For the lower bound, a search \texttt{chkN} on $64$-bit masks takes a state (words still allowed, covered set, forbidden centres), picks the least uncovered word, and branches on the centres that could cover it, forbidding at each branch the siblings already tried.
It prunes when the remaining words cannot cover what is left, at $7$ words per centre.
\texttt{chkN\_\allowbreak sound} proves that a \texttt{true} answer means no completion exists.
A translation argument lets one assume the word $63$ is in the code.
The search from that root is cut at depth $2$ into $38$ states, each refuted by \texttt{decide +kernel} in its own file.
The Lean database of Florath~\cite{florath2026} records $10\le K_2(6,1)\le 12$ at this cell.

\section{The sphere-covering bound}
The library also proves the sphere-covering bound $q^n\le |C|\cdot V$ (\texttt{CoveringA2.\allowbreak sphere\_\allowbreak covering}), its strict form when $V\nmid q^n$, and eight numerical instances \texttt{CoveringChain.\allowbreak SPH\_\allowbreak K*\_\allowbreak lb}, among them $K_7(9,4)\ge 221$.
This was the content of version~0.2 of this note.
The bound had already been formalized by Florath~\cite{florath2026}, and all eight instances are below the best lower bounds in the literature (for $K_7(9,4)$, $264$), so they are not claimed as results.

\section{Other codes, not formalized}
The repository contains seven more codes, each checked by three independent programs outside Lean (Table~\ref{tab:other}).
They are below the upper bounds in the tables we consulted, but none is a Lean theorem yet.
For $q=4$ our literature search was weaker than for $q\ge 5$, and $K_5(10,4)\le 625$ is a linear $[10,4]_5$ code, which should still be compared with the tables of Davydov, Marcugini and Pambianco on linear covering codes.

\begin{table}[ht]
\centering
\begin{tabular}{@{}lrrl@{}}
\toprule
cell & ours & previous & source of previous \\
\midrule
$K_5(7,2)$ & 500 & 525 & Kéri \\
$K_4(10,4)$ & 192 & 208 & Kéri \\
$K_5(9,3)$ & 1250 & 1275 & Kéri \\
$K_5(10,4)$ & 625 & 875 & Kéri \\
$K_5(9,5)$ & 50 & 55 & Kéri \\
$K_5(9,4)$ & 250 & 255 & Kéri \\
$K_7(8,3)$ & 1893 & 2337 & Kéri \\
\bottomrule
\end{tabular}
\caption{Upper bounds verified by computer, not by the kernel.}
\label{tab:other}
\end{table}

\section{What the kernel accepted}
On 1 October 2026, on a virtual machine of type \texttt{e2-highmem-8}, Lean~4.34.1 and Mathlib~v4.34.1:
\texttt{lake build} finished $8942$ jobs in $32$\,s with a peak of $7.7$\,GB;
\texttt{lake build CoveringHeavy} finished $9073$ jobs in $1$\,h\,$41$\,min with up to eight modules in parallel, about $12.4$ CPU-hours in total, at most $8.8$\,GB per process.
\texttt{\#print axioms} on Theorems~\ref{thm:k794}, \ref{thm:ge11} and~\ref{thm:eq12} reports \texttt{propext}, \texttt{Classical.\allowbreak choice} and \texttt{Quot.\allowbreak sound} only.

\section{What is not a theorem}
Novelty is a reading of the literature, not a checked fact.
The seven codes of Table~\ref{tab:other} are not Lean theorems.
That the formal statements say what is intended is a human review: the reader should check the statements of Theorems~\ref{thm:k794}, \ref{thm:ge11}, \ref{thm:eq12} and the definitions of \texttt{Covers} and \texttt{IsK}.
Two files in the tree, \texttt{A5b\_\allowbreak Stress.\allowbreak lean} and \texttt{A6d\_\allowbreak SearchHeavy.\allowbreak lean}, are not imported and are not claimed.

\begin{thebibliography}{9}
\bibitem{cohen1997}
G.~Cohen, I.~Honkala, S.~Litsyn and A.~Lobstein, \emph{Covering Codes}, North-Holland, 1997.
\bibitem{florath2026}
A.~Florath, \emph{Formal foundations and proof-carrying certificates for $q$-ary covering codes in Lean~4}, arXiv:2606.09600, 2026.
\bibitem{gp2025}
D.~Gijswijt and S.~Polak, \emph{Semidefinite lower bounds for covering codes}, arXiv:2504.01932, 2025.
\bibitem{keri}
G.~Kéri, \emph{Tables for bounds on covering codes}, \url{https://old.sztaki.hu/~keri/codes/}, 2011.
\bibitem{lean4}
L.~de~Moura and S.~Ullrich, \emph{The Lean~4 theorem prover and programming language}, CADE~28, LNCS \textbf{12699}, Springer, 2021, 625--635.
\bibitem{marosi2026}
M.~Marosi, \emph{New upper and lower bounds on covering codes $K_q(n,R)$ for alphabets of size $5\le q\le 21$}, arXiv:2608.19872, 2026.
\bibitem{mathlib}
The Mathlib Community, \emph{Mathlib}, \url{https://github.com/leanprover-community/mathlib4}, release v4.34.1.
\bibitem{sk1968}
R.~G. Stanton and J.~G. Kalbfleisch, \emph{Covering problems for dichotomized matchings}, Aequationes Math. \textbf{1} (1968), 94--103.
\end{thebibliography}

\end{document}
